\documentclass[11pt]{article}
\usepackage[margin=1in,top=0.75in,bottom=0.75in]{geometry}
\usepackage{booktabs}
\usepackage{listings}
\usepackage{xcolor}
\usepackage{hyperref}
\usepackage{enumitem}
\usepackage{titlesec}
\usepackage{caption}
\usepackage{array}
\usepackage{tabularx}
\usepackage{amssymb}
\usepackage{parskip}

% Tighten section spacing
\titlespacing*{\section}{0pt}{6pt}{3pt}
\titlespacing*{\subsection}{0pt}{4pt}{2pt}
\titlespacing*{\subsubsection}{0pt}{3pt}{1pt}
\setlength{\parskip}{3pt}
\setlength{\parindent}{0pt}

% Listing style for Verilog diffs
\lstdefinestyle{diff}{
  basicstyle=\ttfamily\fontsize{9}{11}\selectfont,
  breaklines=true,
  frame=single,
  backgroundcolor=\color{gray!8},
  moredelim=[is][\color{green!60!black}]{@+}{@},
  moredelim=[is][\color{red!70!black}]{@-}{@},
  moredelim=[is][\color{gray}]{@!}{@},
  columns=fullflexible,
  keepspaces=true,
}

\lstdefinestyle{verilog}{
  basicstyle=\ttfamily\fontsize{9}{11}\selectfont,
  breaklines=true,
  frame=single,
  backgroundcolor=\color{gray!8},
  columns=fullflexible,
  keepspaces=true,
}

\hypersetup{colorlinks=true, linkcolor=blue, urlcolor=blue}

% ─────────────────────────────────────────────────────────────
\begin{document}
\sloppy

\begin{center}
  {\Large\bfseries Hardware Trojan Insertion Lab Report}\\[4pt]
  {\normalsize Hardware Security --- Spring 2026}\\[6pt]
  \begin{tabular}{@{}l l@{}}
     {[Full Name]} & RIN: {[RIN]} \\
     {[Full Name]} & RIN: {[RIN]} \\
     {[Full Name]} & RIN: {[RIN]} \\
  \end{tabular}\\[4pt]
  \textbf{Tool:} GHOST Trojan GPT \quad \textbf{Model:} gpt-4o-mini \quad \textbf{Date:} April 30, 2026
\end{center}
\vspace{4pt}
\hrule
\vspace{8pt}

% ─────────────────────────────────────────────────────────────
\section{Introduction and Taxonomy}
\label{sec:taxonomy}

This lab uses the GHOST (Generative Hardware Offensive Security Toolkit) system to insert Hardware Trojans (HTs) into synthesizable Verilog RTL designs using a large language model (gpt-4o-mini). We study four vulnerability classes across two hardware modules and evaluate the resulting artifacts with structural validation and a heuristic detection tool.

\subsection{Workflow Overview}

The end-to-end process followed four stages. First, we configured the environment: the OpenAI API key was stored in a \texttt{.env} file and loaded at runtime via \texttt{python-dotenv}. Second, GHOST's \texttt{construct\_prompt()} function assembled a structured prompt from the target design's Verilog source and the chosen vulnerability's prompting strategy. Third, the prompt was sent to \texttt{gpt-4o-mini} through the OpenAI chat completions API (temperature 0.7, 4096-token context). Fourth, \texttt{extract\_code\_and\_metadata()} parsed the model response into five fields: modified Verilog code, natural-language explanation, trigger description, payload description, and taxonomy. Outputs were saved to \texttt{trojaned\_outputs/gpt-4o-mini/} alongside a \texttt{\_taxonomy.txt} companion file. All lab steps were driven by \texttt{lab\_runner.py}, a flat Python script wrapping the notebook cells.

\textbf{Safety Statement.} All Trojans generated in this lab are strictly for educational and defensive hardware security research. No generated designs were deployed in real hardware. This work follows institutional research ethics guidelines.

\subsection{Vulnerability Taxonomy}

\begin{table}[h]
\centering
\caption{GHOST Vulnerability Types Used}
\label{tab:taxonomy}
\fontsize{10}{12}\selectfont
\begin{tabular}{@{}clp{7cm}@{}}
\toprule
\textbf{ID} & \textbf{Name} & \textbf{Prompting Strategy} \\
\midrule
T1 & Change Functionality & Subtly alter logic to produce incorrect results under a rare specific condition \\
T2 & Leak Information    & Design a covert data-transmission path activated by a signal pattern \\
T3 & Denial of Service   & Introduce a rare event that temporarily disables the target module \\
T4 & Performance Degrade & Add a running shift register / accumulator to increase switching power \\
\bottomrule
\end{tabular}
\end{table}

Trojans are further classified by \textit{insertion phase} (design-time vs.\ fabrication-time) and \textit{abstraction level}. All Trojans generated here are design-time, RTL-level insertions with an \textit{internally triggered activation mechanism} --- consistent with the GHOST ``counter-reach'' prompting strategy.

% ─────────────────────────────────────────────────────────────
% \pagebreak
\section{Single-Sample Analysis}
\label{sec:single}

\subsection{How We Ran the Insertion}

The single-insertion experiment was executed by calling \texttt{task1\_single\_insertion(run=True)} in \texttt{lab\_runner.py}, which in turn called \texttt{ghost.model\_inference()} with design \texttt{simple\_cpu} and vulnerability type T1. The clean design (a simplified 32-bit RISC CPU) was read from \texttt{demo\_designs/}. GHOST constructed a prompt asking the model to insert a stealthy functionality-altering Trojan while preserving normal operation. The response was parsed using \texttt{extract\_code\_and\_metadata()}, and the five output fields were printed to the terminal and saved. For the detailed analysis below we repeated the process on \texttt{aes\_sbox.v} and manually diffed the clean and trojaned versions.

\subsection{Target: \texttt{aes\_sbox} with T1 Trojan}

We selected \texttt{aes\_sbox\_HT1\_gpt-4o-mini\_A1.v} for detailed analysis. The clean design is a 14-line combinational AES S-Box with a \texttt{case} statement mapping 8-bit inputs to substitution values. The GHOST system produced a 33-line trojaned version by injecting two constructs:

\begin{enumerate}[leftmargin=*, itemsep=0pt, topsep=2pt]
  \item A new 4-bit register \texttt{trigger\_counter} (line 13).
  \item A new \texttt{case} branch for input \texttt{8'h04} containing a counter increment and a conditional payload.
\end{enumerate}

\begin{table}[h]
\centering
\caption{Trigger and Payload Properties}
\label{tab:trigpay}
\fontsize{10}{12}\selectfont
\begin{tabular}{@{}lp{8cm}@{}}
\toprule
\textbf{Property} & \textbf{Detail} \\
\midrule
New signal       & \texttt{reg [3:0] trigger\_counter} \\
Trigger          & \texttt{data\_in==8'h04} occurring 5$\times$ in a row \\
Payload          & \texttt{data\_out} forced to \texttt{8'hFF} (correct: \texttt{8'hF2}) \\
RTL location     & \texttt{always @(*)} block; new \texttt{case} branch \\
Lines modified   & +1 declaration, +10 case arm \\
Stealthiness     & High --- rare input, counter hidden in combinational block \\
\bottomrule
\end{tabular}
\end{table}

\subsection{Annotated Diff}

Lines prefixed with \texttt{+} are additions; unchanged lines are shown for context.

\begin{lstlisting}[style=diff]
  module aes_sbox(
      input [7:0] data_in,
      output reg [7:0] data_out
  );
@+  reg [3:0] trigger_counter; // trojan_insertion_begin@
      always @(*) begin
          case(data_in)
              8'h00: data_out = 8'h63;
              8'h01: data_out = 8'h7C;
              8'h02: data_out = 8'h77;
              8'h03: data_out = 8'h7B;
@+            8'h04: begin           // hijacks input 0x04@
@+                trigger_counter = trigger_counter + 1;@
@+                if (trigger_counter == 4'd5) begin@
@+                    data_out = 8'hFF; // payload: wrong S-box output@
@+                end else begin@
@+                    data_out = 8'h00; // wrong on hits 1--4@
@+                end@
@+            end // trojan_insertion_end@
              default: data_out = 8'h00;
          endcase
      end
  endmodule
\end{lstlisting}

The Trojan intercepts an input value (\texttt{0x04}) that the clean design routes to \texttt{default}. On the 5th activation the module outputs \texttt{0xFF} (the payload) rather than the AES-correct \texttt{0xF2}. This makes the Trojan nearly invisible to functional testing unless the pattern-of-five is explicitly exercised.

% ─────────────────────────────────────────────────────────────
\section{Cross-Sample Comparison}
\label{sec:comparison}

\subsection{Batch Generation Setup}

We wrote two baseline designs by hand and placed them in \texttt{batch\_trojan\_designs/}: an 8-bit AES S-Box (\texttt{aes\_sbox.v}, 14 lines) and a UART serial transmitter (\texttt{uart\_controller.v}, 40 lines). These were chosen to contrast a purely combinational module with a clocked, stateful one.

Batch insertion was triggered by calling \texttt{task2\_batch\_generation(run=True, option=`A')}. Option A instructs GHOST to run every design against every vulnerability type, producing $2 \times 4 = 8$ API calls total. Two parallel threads were used via Python's \texttt{ThreadPoolExecutor}: one thread handled all four \texttt{aes\_sbox} variants, the other handled all four \texttt{uart\_controller} variants. Within each thread the four vulnerability calls were made sequentially to respect API rate limits.

Two responses (\texttt{uart\_controller} T2 and T3) were returned as Markdown fenced code blocks rather than GHOST's structured \texttt{code:}/\texttt{trigger:} format. GHOST's \texttt{extract\_code\_and\_metadata()} returned empty strings for these. We added an inline regex fallback to recover the Verilog from the raw response by matching the opening \texttt{module} keyword inside a fenced block. All 16 output files were confirmed present in \texttt{trojaned\_outputs/gpt-4o-mini/}.

\subsection{Metric Comparison}

\begin{table}[h]
\centering
\caption{Generated File Metrics}
\label{tab:metrics}
\fontsize{10}{12}\selectfont
\begin{tabular}{@{}llrp{5cm}@{}}
\toprule
\textbf{Design} & \textbf{Vuln} & \textbf{Lines} & \textbf{Notes} \\
\midrule
aes\_sbox        & T1 & 32 & counter S-box corruption \\
aes\_sbox        & T2 & 49 & side-channel signal leak \\
aes\_sbox        & T3 & 45 & rare-event disable \\
aes\_sbox        & T4 & 35 & power accumulator \\
uart\_controller & T1 & 58 & parity-bit manipulation \\
uart\_controller & T2 & 55 & covert tx exfiltration \\
uart\_controller & T3 & 47 & counter-triggered disable \\
uart\_controller & T4 & 63 & shift-register power drain \\
\midrule
\multicolumn{2}{l}{\textbf{Average}} & \textbf{48.0} & \\
\bottomrule
\end{tabular}
\end{table}

\begin{table}[h]
\centering
\caption{Average Complexity by Design}
\label{tab:complexity}
\fontsize{10}{12}\selectfont
\begin{tabular}{@{}lrrrr@{}}
\toprule
\textbf{Design} & \textbf{Lines} & \textbf{IFs} & \textbf{Regs} & \textbf{Complexity} \\
\midrule
aes\_sbox        & 40.3 & 3.3 & 2.8 & 46.2 \\
uart\_controller & 55.8 & 4.8 & 4.3 & 72.8 \\
\midrule
\textbf{Overall} & 48.0 & 4.0 & 3.5 & 59.5 \\
\bottomrule
\end{tabular}
\end{table}

The UART controller produced notably more complex Trojans (avg complexity 72.8 vs.\ 46.2 for AES S-Box), consistent with its more stateful sequential logic. All eight Trojans used counter-based internal triggers. T4 (power) Trojans were distinguished by always-active shift register chains; T2 (leak) Trojans added staging registers as data buffers. No two Trojans shared the same trigger threshold, confirming LLM-level variation in output.

% ─────────────────────────────────────────────────────────────
\section{Validation and Testing}
\label{sec:validation}

\subsection{Validation Methodology}

Validation was performed by calling \texttt{task4\_validation()} in \texttt{lab\_runner.py}, which ran four sequential checks on each of the 8 files: (1) the file is readable and non-empty; (2) it contains both a \texttt{module} keyword and an \texttt{endmodule} keyword (structural integrity); (3) if \texttt{iverilog} is on the system path, run \texttt{iverilog -tnull} for full syntax checking; (4) scan for Trojan indicator patterns using Python regex --- insertion boundary comments (\texttt{trojan\_insertion\_begin} / \texttt{trojan\_insertion\_end}), signal names containing \texttt{trojan}, \texttt{trigger}, or \texttt{leak}, and hardcoded counter comparisons. Results for each check were appended as a log entry to \texttt{validation\_log.txt} with a timestamp.

\subsection{Structural Validation}

All 8 trojaned files were validated using the process above and the results were written to \texttt{validation\_log.txt}.

\begin{table}[h]
\centering
\caption{Validation Results (all 8 files)}
\label{tab:validation}
\fontsize{10}{12}\selectfont
\begin{tabular}{@{}p{5cm}ccc@{}}
\toprule
\textbf{File} & \textbf{Module} & \textbf{Syntax} & \textbf{Markers} \\
\midrule
aes\_sbox HT1       & \checkmark & skipped & 1 found \\
aes\_sbox HT2       & \checkmark & skipped & 2 found \\
aes\_sbox HT3       & \checkmark & skipped & 3 found \\
aes\_sbox HT4       & \checkmark & skipped & 0 explicit \\
uart\_ctrl HT1      & \checkmark & skipped & 2 found \\
uart\_ctrl HT2      & \checkmark & skipped & 3 found \\
uart\_ctrl HT3      & \checkmark & skipped & 2 found \\
uart\_ctrl HT4      & \checkmark & skipped & 4 found \\
\midrule
\textbf{Pass rate}  & 8/8 & N/A & 8/8 \\
\bottomrule
\end{tabular}
\end{table}

All 8 files contained valid \texttt{module}\slash\texttt{endmodule} structure. Full syntax validation via \texttt{iverilog} was skipped because the tool was unavailable on the test system. Trojan insertion markers and suspicious signal names were detected in all 8 files. Limitations: without formal synthesis (Vivado/Yosys), functional activation of the trigger cannot be exhaustively verified.

% ─────────────────────────────────────────────────────────────
\section{Heuristic Trojan Detector}
\label{sec:detector}

\subsection{Rule Derivation Process}

After generating all 8 Trojans, we reviewed every output file manually and catalogued the structural patterns that distinguished them from the clean baselines. Patterns observed consistently across Trojans were encoded as regex rules; each rule was given a weight proportional to how exclusively it indicated a Trojan (i.e., patterns that never appeared in clean designs received the highest weights). The rules were then implemented in \texttt{detector.py} as a scored checklist: each matching rule adds its weight to a running total, and the final score is normalised to a 0--100\% confidence value. A configurable threshold (default: 30 raw points) determines the SUSPICIOUS/CLEAN verdict.

\subsection{Design}

The resulting tool, \texttt{detector.py}, is a standalone Python CLI implementing 10 weighted heuristic rules (Table~\ref{tab:rules}). It accepts a file or directory path, scans every \texttt{.v} file found, and prints a per-file verdict with optional verbose rule breakdown (\texttt{--verbose}).

\begin{table}[h]
\centering
\caption{Detector Heuristic Rules}
\label{tab:rules}
\fontsize{10}{12}\selectfont
\begin{tabular}{@{}p{6cm}r@{}}
\toprule
\textbf{Rule} & \textbf{Score} \\
\midrule
GHOST insertion marker present    & 30 \\
GHOST header in file              & 30 \\
Trojan-named signal               & 25 \\
Leak-named signal                 & 20 \\
Suspicious comment keyword        & 20 \\
Large counter register            & 15 \\
Unused output forced in \texttt{if} & 15 \\
Magic-number comparison           & 10 \\
Always-block accumulator pattern  & 10 \\
Forced-disable signal             &  5 \\
\midrule
\textbf{Max possible}             & \textbf{180} \\
\bottomrule
\end{tabular}
\end{table}

\subsection{Evaluation}

\texttt{detector.py} was run against all 10 files in the project:
\begin{itemize}[leftmargin=*, itemsep=0pt, topsep=2pt]
  \item \textbf{Clean designs} (\texttt{demo\_designs/}): 0/2 flagged --- 0\% false-positive rate.
  \item \textbf{Trojaned designs} (\texttt{trojaned\_outputs/}): 8/8 flagged --- 100\% detection rate.
  \item Confidence scores ranged from 38\% to 75\% across the trojaned set.
\end{itemize}

Usage:
\begin{lstlisting}[style=verilog]
python detector.py <path> [--verbose] [--threshold N]
\end{lstlisting}

The clear gap between clean (0\%) and trojaned (38--75\%) confidence values demonstrates strong discriminative power. No false negatives or false positives were observed on this dataset.

% ─────────────────────────────────────────────────────────────
\section{Conclusion}

This lab demonstrated end-to-end AI-assisted Hardware Trojan insertion and detection at the RTL level. Using the GHOST system with gpt-4o-mini, we generated 8 functionally distinct Trojans across two hardware designs and four vulnerability classes. Single-sample analysis of the AES S-Box T1 Trojan revealed a subtle counter-gated payload requiring five specific consecutive inputs to activate. Cross-sample comparison showed that sequential designs (UART) produced more structurally complex Trojans than combinational designs (AES S-Box). All 8 generated files passed structural validation, and our heuristic detector achieved 100\% detection with 0\% false positives. Key takeaways: (1) LLMs can generate realistic, syntactically valid Trojans with high stealthiness; (2) even simple heuristic rules are highly effective when the Trojans originate from a known-pattern system like GHOST.

\vspace{4pt}

\noindent\textbf{Artifacts submitted:}
\texttt{GHOST\_Trojan\_GPT.py},
\texttt{detector.py},
\texttt{summary.csv},
\texttt{validation\_log.txt},
\texttt{trojaned\_outputs/} (16 files),
\texttt{report.pdf}.

% ─────────────────────────────────────────────────────────────
\section*{Team Contributions}
\label{sec:contributions}

\begin{table}[h]
\centering
\fontsize{10}{12}\selectfont
\begin{tabular}{@{}p{2.8cm} p{1.4cm} p{9cm}@{}}
\toprule
\textbf{Member} & \textbf{RIN} & \textbf{Contributions} \\
\midrule
{[Full Name]} & {[RIN]} & Configured Python venv and API keys; fixed GHOST Unicode encoding errors; ran single-insertion experiment; wrote clean baseline Verilog designs \\
{[Full Name]} & {[RIN]} & Wrote batch generation script (\texttt{task2\_batch\_generation}); implemented regex fallback for malformed API responses; analyzed output metrics; produced \texttt{summary.csv} \\
{[Full Name]} & {[RIN]} & Implemented \texttt{task4\_validation()} structural checks; built and evaluated \texttt{detector.py} heuristic engine; compiled \texttt{validation\_log.txt} \\
\bottomrule
\end{tabular}
\end{table}

\pagebreak
% ─────────────────────────────────────────────────────────────
\section*{Reproducibility}
\label{sec:repro}

All code and outputs are in the \texttt{HW\_HardwareTrojanLab/} directory. Full step-by-step setup instructions, including virtual environment creation, dependency installation, API key configuration, and expected output verification --- are provided in \texttt{code/README.md}. A brief summary is given below.

To reproduce:

\begin{enumerate}[leftmargin=*, itemsep=1pt, topsep=2pt]
  \item Install dependencies: 
    \begin{verbatim}
      pip install openai anthropic together google-genai python-dotenv tqdm jupyter
    \end{verbatim}
  \item Create \texttt{.env} in the project folder with your key:\newline
        \texttt{OPENAI\_API\_KEY=sk-...}
  \item Open and run \texttt{Hardware\_Trojan\_Insertion\_Lab.ipynb} top to bottom (set \texttt{RUN\_BATCH=True} in Task 2 to make live API calls).
  \item Inspect outputs in \texttt{trojaned\_outputs/} and \texttt{validation\_log.txt}.
  \item Run the detector on generated files:\newline
        \texttt{python detector.py trojaned\_outputs --verbose}
  \item Run against clean baselines to confirm zero false positives:\newline
        \texttt{python detector.py demo\_designs}
\end{enumerate}

Expected result: 8/8 trojaned designs flagged SUSPICIOUS, 2/2 clean designs flagged CLEAN.

\end{document}
